\section{Discussion of  results}

\subsection{Rest-frame density and $Z$ factor}



An important object  is the rest-frame density
${\cal M} (0,z_3^2)$. It is produced by data at $P=0$.
The results for  its imaginary part are compatible with zero,
as required.
The real part, shown in Fig. \ref{M0},   is a symmetric function of $z_3$,
and has a clearly visible linear component in its  fall-off with  $|z_3|$ for  small and
middle values  of $|z_3|$.  In fact, a  linear exponential factor
$Z(z_3^2) \sim  e^{-c |z_3|/a}$ is expected as a manifestation
of the nonperturbative effects generated by  the straight-line gauge  link.


\begin{figure}[t]
  \centerline{\includegraphics[width=3.7in]{images/M0NormK2.pdf}}
    \caption{Real part of the rest-frame  density  ${\cal M} (0,z_3^2)$
    \label{M0}}
    \end{figure}


\subsection{Reduced Ioffe-time distributions}

 In   Fig. \ref{realz},  we plot the results for
the real part of the  ratio  ${\cal M} (Pz_3, z_3^2)/{\cal M} (0,z_3^2)$ as a function of $z_3$
taken at six fixed values of the momentum $P$.
One can see that  all the curves  have  a Gaussian-like shape.
Thus, the $Z(z_3^2)$  link renormalization factor has been canceled in the ratio,
as expected.



Furthermore, the curves  look  similar  to each other, {differing only by a decreasing  width  with $P$}.
In  Fig. \ref{realc} , we plot   the  same data, but change the axis to $\nu=Pz_3$.
As one can see, now the  data practically fall on the same curve.
For  the  imaginary part, the situation is similar.

 This phenomenon corresponds to factorization of the $x$- and $k_\perp$-dependence
    for the   soft TMD ${\cal F} (x, k_\perp^2)$, as discussed in previous sections.


\begin{figure}[t]
    \centerline{\includegraphics[width=4.0in]{images/ReMz.pdf} }
    \caption{Real  part of the reduced  distribution ${\mathfrak  M} (Pz_3,z_3^2)$ plotted as
    a  function of
    $z_3$. Here, $P= 2 \pi p/L$.
    \label{realz}}
    \end{figure}


        \begin{figure}[t]
    \centerline{\includegraphics[width=3.8in]{images/RealC.pdf}}
    \caption{Real  part of ${\mathfrak  M} (\nu, z_3^2)$  plotted as a function
    of $\nu =Pz_3$ and
    compared to the curve given by Eqs. (\relax {\ref{MC}}), (\relax {\ref{qV}}).
        \label{realc}}
    \end{figure}

      \begin{figure}[t]
    \centerline{\includegraphics[width=3.8in]{images/pheno.pdf} }
    \caption{Valence  distribution $q_v(x) $ as given  by \mbox{Eq. (\relax {\ref{qV}})}
    compared with the  $Q^2=1$ GeV$^2$  NNLO global fits NNPDF31\_nnlo\_pch\_as\_0118\_mc\_164 \protect{\cite{Ball:2017nwa}  } and
  MSTW2008nnlo68cl\_nf4   \protect{\cite{Martin:2009iq}}; and the NLO global fit CJ15nlo  \protect{\cite{Accardi:2016qay}},
  all extracted using the LHAPD6  library  \protect{\cite{Buckley:2014ana}}.
  The bands around the global fits indicate their experimental and systematic uncertainties.
        \label{MSTW}}
    \end{figure}


    \subsection{Quark-antiquark decomposition}

    The real part of the Ioffe-time distribution is obtained from  the cosine Fourier transform
  \begin{align}
{\mathfrak  M}_R (\nu) \equiv   & \int_0^1 dx \,
\cos (\nu x)  \, q_v  (x)
\label{MC}
     \end{align}
of the function  $q_v(x)$  given by the
difference  $q_v(x) = q(x) -\bar q(x)$  of
quark and antiquark   distributions.
 In our case,   $q$ is  $u-d$ and
    $\bar q = \bar u - \bar d$.
    The $x$-integral of $u-\bar u$ equals to the number of $u$-quarks
    in the proton, which is 2, while
    the  $x$-integral of $d- \bar d$  equals 1. Thus, the $x$-integral of $q_v (x)$ should be equal to 1.



We found that our  data for the real part are well described  if one
chooses       the function
     \begin{align}
    q_v(x) =  \frac{315}{32} \sqrt{x}  (1-x)^{3}\,  ,
    \label{qV}
    \end{align}
whose $x$-integral is normalized to 1.
To get it, we formed cosine Fourier transforms ${\cal M} (\nu;a,b)$ of the
normalized $x^a(1-x)^b$-type  functions
and found the parameters $a,b$  by  fitting    our data.
The comparison of the  data with the  curve based on
Eqs. (\relax {\ref{MC}}), (\relax {\ref{qV}})  is shown in Fig.  \ref{realc}.

While all the data points were used in the fit, the latter  is clearly dominated by  %
the points with the smaller values of   Re\ ${\mathfrak  M} (\nu, z_3^2)$.   %
For $\nu <10$,  the data points  lying  above the curve, correspond to  %
values of $z_3=3a$ to $5a$. As we will see later, they  reflect the perturbative evolution:  %
\mbox{Re \ ${\mathfrak  M} (\nu, z_3)$}  increases when $z_3$ decreases.  %
In this context, the overall  curve (\ref{qV}) corresponds to PDF  ``at low normalization point'', %
i.e.,  in  the region, where the perturbative evolution stops. %

In general,  it is more  appropriate to fit  Re  ${\mathfrak  M} (\nu, z_3^2)$ as a function of two variables,  %
$\nu$ and $ z_3$, even though the dependence on $z_3$  is rather weak and  %
noticeable  just for a few points. Since  we made a fit of  Re  ${\mathfrak  M} (\nu, z_3)$ %
as a function of just one variable $\nu$, there are points that visibly deviate
from the curve, but we do not think that  %
 it  makes    sense to translate the evolution $z_3$-dependence %
of small-$z_3$  points into    %
an error band to our curve in \mbox{Fig.  \ref{realc}. }  %
In Section 5.4, we evolve the data points to a common reference scale $z_0=2a %%%%%%
$ and show the error band for the results obtained in this way.  %%%%%%%%%%

We realize   that our lattice setup  is rather crude (quenched approximation, very large pion mass), %%%%%%
and for this reason  we do not attempt  to perform a thorough comparison of our results %%%%%%%%%
with experimental data. Still,  we think that some  kind of comparison is rather useful as  an illustration.  %%%%%%%%%

Thus, we %%%%%%%%%
compare our $q_v (x) $ with three  global fits
for the difference \mbox{$u_v(x) -d_v(x)$}  of the valence distributions,  see Fig. \ref{MSTW}.
These global  fits curves correspond to   $\mu=1$ GeV scale, while our  %%%%%%%%%
``low normalization point'' curve corresponds to $\mu \lesssim 0.3$ GeV.  %%%%%%%%
Still, one can see that our curve is not very far from the %%%%%%%%%%
NNPDF31  \cite{Ball:2017nwa}  NNLO fit down  to $x=0.1$
 and from the   MSTW  \cite{Martin:2009iq} NNLO  fit  down  to $x=0.05$.
 We also show the NLO fit CJ15  \cite{Accardi:2016qay}.




Since the areas  under each  curve  are  equal to 1, our curve
compensates the strong deficiency  in the $x<0.1$ region
by exceeding the NNLO curves  at $x>0.1$ values.
In other words, if our curve would better describe data in the $x<0.1$ region,
it would necessarily be smaller in the $x>0.1$ region.


      \begin{figure}[t]
    \centerline{\includegraphics[width=3.9in]{images/ImC.pdf} }
    \caption{Imaginary   part of ${\mathfrak  M} (\nu, z_3^2)$
    compared to the curve   $  {\mathfrak  M}_I^v (\nu)$  based on  $\bar q (x)=0$.
        \label{imc}}
    \end{figure}


      \begin{figure}[t]
    \centerline{ \includegraphics[width=3.9in]{images/ImCA.pdf}}
    \caption{Imaginary   part of ${\mathfrak  M} (\nu, z_3^2)$
    compared to the curve based on  $\bar q(x)$ given by Eq. (\relax {\ref{qA}}).
        \label{imca}}
    \end{figure}



The sine Fourier transform
      \begin{align}
  {\mathfrak  M}_I (\nu)  \equiv
& \int_0^1 dx \,
\sin (\nu x) \, q_+  (x)
     \end{align}
   is  built from  the function $q_+ (x)= q(x) + \bar q (x)$,
   which may be also represented as \mbox{$q_+ (x)= q_v(x) + 2 \bar q (x)$. }
   If we neglect the antiquark contribution and use \mbox{$q_+  (x) = q_v(x)$,}
   we get the curve shown in   Fig. \ref{imc}
    (call   it $  {\mathfrak  M}_I^v (\nu)$). The agreement with the data
   is strongly improved if we use a non-vanishing antiquark contribution, namely
        \begin{align}
    \bar q(x) =\bar u (x) - \bar d(x) =  0.07 \left [ 20  \,  x (1-x)^{3} \right ] \,  ,
    \label{qA}
    \end{align}
    see Fig. \ref{imca}.
    This function was obtained by fitting the data for %%
    the difference ${\rm Im} \ {\mathfrak M} (\nu, z_3^2) - {\mathfrak M}_I^v(\nu)$ %%%%%%%
    by sine Fourier transforms of $A x^a (1-x)^b$ functions. %%%%%%%%
     This result corresponds to
          \begin{align}
 \int_0^1 dx \, [\bar u (x) - \bar d(x)] =  0.07 \ .
     \end{align}
     The combined distribution
         \begin{align}
q(x)& = u(x)-d(x)    \nn & =  [q_v (x) +\bar q (x)  ]\, \theta(x>0)
-
\bar q (-x)  \, \theta(x<0)
\label{q}
     \end{align}
 defined on the $-1\leq x \leq 1$ interval
     is shown in Fig. \ref{umind}.


          \begin{figure}[t]
    \centerline{\includegraphics[width=4.0in]{images/umindkey.pdf} }
    \caption{Overall distribution $q(x) $ as defined by Eq. (\relax {\ref{q}}).
        \label{umind}}
    \end{figure}


        \subsection{Evolution}


    While an overall agreement of the data
    with a \mbox{$z_3$-independent} curve  looks satisfactory,
    one can easily notice  a residual \mbox{$z_3$-dependence}  in the data.
    It  is  especially visible when, for a particular $\nu$,    there  are several  data points
  corresponding to different values of $z_3$.
It is interesting to check if this dependence corresponds to perturbative evolution.




To begin with,  the  evolution of the real part should  lead
to its  {\it decrease} when $z_3^2$ increases. On the other hand, as pointed out at the end of Section 3,
 the function  ${{\rm Im} \, \cal M} ( \nu, z_3^2)$
{\it increases}
 when $z_3^2$ {\it increases}  as long as  \mbox{$\nu \lesssim 5.5$.}
 Our data follow these patterns.

As we discussed, the evolution corresponds to $\ln z_3^2$  singularities of the Ioffe-time distributions
for small $z_3^2$.  Thus, a natural idea is  to check if the data  corresponding to
small  $z'_3$ and $z_3$  may be related by
  \begin{align}
{\mathfrak M} (\nu, {z'}_3^2) {=}
{\mathfrak M} (\nu, z_3^2) \
 -\frac23  \frac{\alpha_s}{\pi}   \ln ({z'_3}^2/z_3^2) B \otimes {\mathfrak M}\,  (\nu, z_3^2)
 \label{Prop}
 \end{align}
 for some value of $\alpha_s$.
Here $B$ is the evolution kernel  (\ref{Bu}).
In our case,
    \begin{align}
B \otimes {\mathfrak M}\,  (\nu)   & =
\int_0^1  du \,  \frac{1+u^2}{1- u} \,    [{\mathfrak M} ( \nu)- {\mathfrak  M} (u \nu)  ]
 .
\label{plusM}
 \end{align}





More specifically, we fix  the point  $z'_3$ at the   value \mbox{$z_0=2a$}   corresponding,
at the leading logarithm level,  %%%%%%%%
 to the
$\overline {\rm MS}$-scheme scale \mbox{$\mu_0 = 1$  GeV }  and build the function
   \begin{align}
\widetilde{\mathfrak M} (\nu, z_0^2) {\equiv }
{\mathfrak M} (\nu, z_3^2) \
 -\frac23  \frac{\alpha_s}{\pi}   \ln (z_0^2/z_3^2) B \otimes {\mathfrak M}\,  (\nu, z_3^2)
 \label{EvoM}
 \end{align}
 from the data points for ${\mathfrak M}\,  (\nu, z_3^2)$
 using various values for $\alpha_s$.

 Since the perturbative evolution is expected for small $z_3$, we include in this analysis
 the data with $z_3$ up to 4
lattice spacings, which corresponds to energy scales \mbox{$\mu = 2,     1,     0.7$  and     0.5  GeV.}


                \begin{figure}[t]
    \centerline{ \includegraphics[width=3.9in]{images/EvoRe0.pdf}}
    \caption{Real    part of ${\mathfrak  M} (\nu, z_3^2)$
    for $z_3/a =1, 2, 3,$ and 4.
        \label{ER0}}
    \end{figure}

      \begin{figure}[t]
    \centerline{ \includegraphics[width=3.8in]{images/EvoReEv.pdf}}
    \caption{Evolved data points for the real part.
        \label{ERE}}
    \end{figure}



For the real part, these data points are shown in  \mbox{Fig. \ref{ER0}. }
 As one can see, there is a visible scatter of the data points.
    Using $\alpha_s/\pi =0.1$, we calculate the
    ``evolved'' data points corresponding to  the function $\widetilde{\mathfrak M} (\nu, z_0^2) $.
    The results are shown in Fig. \ref{ERE}.  The evolved data points  are now very close
    to a universal  curve.



    \begin{figure}[t]
    \centerline{ \includegraphics[width=3.8in]{images/EvoIm0.pdf}}
    \caption{Imaginary     part of ${\mathfrak  M} (\nu, z_3^2)$
    for $z_3/a =1, 2, 3,$ and 4.
        \label{EI0}}
    \end{figure}

           \begin{figure}[t]
    \centerline{ \includegraphics[width=3.8in]{images/EvoImEv.pdf}}
    \caption{Evolved data points for the imaginary part.
        \label{EIE}}
    \end{figure}


        \begin{figure}[t]
    \centerline{ \includegraphics[width=4.0in]{images/Revol10.pdf}}
    \caption{Data points for Re ${\mathfrak M}\,  (\nu, z_3^2)$ with $z_3 \leq 10a$
     evolved to $z_0=2a$ as described in the text.
        \label{evpoi}}
    \end{figure}


        \begin{figure}[t]
    \includegraphics[width=4.2in]{images/evol4.pdf}
    \caption{
    Curve for $u_v(x) - d_v(x)$ built from the evolved data
   shown in Fig. \ref{evpoi}, and     treated as corresponding to
     the  \mbox{$\mu^2 =1$  GeV$^2$}   scale;
then  evolved to the reference point  $\mu^2 =4$  GeV$^2$  of the  global fits.
        \label{evolved}}
    \end{figure}



 In  Fig. \ref{EI0},    we show the initial data points for the imaginary  part. The evolved data points
 constructed using the same  $\alpha_s/\pi =0.1$ value  are shown  in Fig. \ref{EIE}.
        Again, they are close to a universal curve.
        This analysis indicates that the residual $z_3^2$-dependence
         of ${\mathfrak M}\,  (\nu, z_3^2)$ at fixed $\nu$ is compatible with the expected logarithmic evolution at small $z_3^2$.
Clearly this is an important feature of our calculation which needs to be further studied as it will play an essential role in reliable extraction of renormalized PDFs from this type of lattice calculations.


With a smaller lattice spacing, the use of perturbative evolution may be justified  %%%%%%%%
in a wider region of $\nu$. While our data extend to rather large %%%%%%%%
separations $\sim 1$ fm, we find it instructive  to use them as an example to   %%%%%%%%%
 illustrate the trends generated  by %%%%%%%
the perturbative evolution.    %%%%%%%%

To this end, we applied the leading logarithm formula %%%%%%%%
(\ref{EvoM}) with $z_0=2a$ and $\alpha_s/\pi=0.1$ to our data points with $z_3 \leq 6a$.  %%%%%%
Assuming that evolution stops for $z_3 \gtrsim 6a$ (as indicated by our data), %%%%%%%%%%
the data points with \mbox{$7a \leq z_3 \leq 10a$}   %%%%%%%%%%
 were evolved to $z_0$ using Eq. (\ref{EvoM}) with $z_3=6a$. %%%%%%%%%%
The  data points evolved in this way are shown in Fig. \ref{evpoi}.    %%%%%%%%%%


Fitting  the evolved  points by  cosine Fourier transforms ${\cal M} (\nu;a,b)$ of the %%%%%%%%%%
normalized $N(a,b) x^a(1-x)^b$-type  functions,  %%%%%%%%%%%
we  found that they may be described if one takes $a = 0.36(6)$  and $b=3.95(22)$.  %%%%%%%%%%%
Treating $z_0=2a$ as the $\overline{\rm MS}$ scale $\mu=1$ GeV, one can further evolve  %%%%%%%%%%%
the curve to the standard reference scale $\mu^2 = 4$ GeV$^2$ of the global fits, see Fig. \ref{evolved}.    %%%%%%%%%%%
Comparing with Fig. \ref{MSTW}, we see that the perturbative evolution shifts   %%%%%%%%%%%
our curves,  moving them  closer to the global fits.   %%%%%%%%%%%
